\documentclass[main.tex]{subfiles}


\begin{document}

% \AddToShipoutPictureFG{%
% \begin{tikzpicture}[overlay,remember picture]
% \path (current page.south west) -- (current page.north east)
% node[midway,scale=1,color=gray,sloped,opacity=0.4] {\textnormal{Кравченко Игорь Игоревич\hspace{1cm} +79010144910\hspace{1cm} super.antig@yandex.ru}};
% \end{tikzpicture}
% }


% %from pagenumber to up
% \phantomsection 
% \hypertarget{toc}{}

 

 

% номер секции вручную
\setcounter{section}{17
}
\addtocounter{section}{-1} 

\section{Второй и третий законы Ньютона}


Если результирующая сила, действующая на тело, \emph{не равна нулю}, то в ИСО\footnote{Инерциальная система отсчета.} тело движется с ускорением. На рис.\,\ref{fig:p33} показан опыт с толканием тележки.

\begin{figure}[H]
    \centering

    \begin{tikzpicture}[font=\small,            line cap=round,                             >={Stealth[inset=0pt,round,length=3mm, angle'=20]},vec/.style={>={Stealth[inset=0pt,round,length=3.5mm, angle'=20]}}]


\filldraw[lightgray,semitransparent] (0,0) rectangle (12,-0.3);
\draw (0,0) -- (12,0);

\begin{scope}[xshift=4cm]
    
%%motion1
\filldraw[lightgray,draw=black] (0,0.5) rectangle (4,1);
\node [above,black] at (2,1) {Т};
\draw[->,Green,thick,vec] (3,1.2) --++ (1,0) node[black,above,midway] {$\vec a_\text{т}$};

%wheels
\draw[black,fill=black]
    (1,.3) circle(.3cm)
    (3,.3) circle(.3cm);
\draw[fill=gray]
    (1,.3) circle(.2cm)
    (3,.3) circle(.2cm);
\end{scope}



%%%%%%%%%%%%%%%%%%%%%%%%%%%%%%%%%%%%%%
% PERSON
  \def\W{3.6}  % ground width
  \def\D{0.2}  % ground depth
  \def\H{1.9}  % human height
  \def\h{0.9}  % mass height
  \def\w{1.0}  % mass width
  \def\mx{0.10*\W} % mass x coordinate

\begin{scope}[xshift=3.6cm]
    
  \coordinate (H) at (0,0.75*\H);
  \draw[thick,fill=white] (H) circle (0.3) node[xshift=-0.3cm,left] {$\text{Н}$};
  \draw[thick] (H)++(-120:0.3) coordinate (N) to[out=-115,in=40]++ (-130:0.36*\H) coordinate (P);
  \draw[thick] (N)++(-115:0.03) to[out=-60,in=170] (.4,\h);
  \draw[thick] (N)++(-115:0.03) to[out=-70,in=170] (.4,0.92*\h);
  \draw[thick] (P) to[out=-120,in=30,looseness=1.3] (-0.35*\W,0);
  \draw[thick] (P) to[out=-50,in=30,looseness=1.8] (-0.22*\W,0);

  %F
  % \node at (P) [circle,fill=red,inner sep=0pt,minimum size=1mm,label=below:] {};
  % \draw[->,red,thick,vec] (P) --++ (-3,0) node[black,above,midway] {$\vec F_\text{тележки}$};
      
\end{scope}  

%vec
\node at (6,.75) [circle,fill=red,inner sep=0pt,minimum size=1mm,label=below:] {};
\draw[->,red,thick,vec] (6,.75) --++ (3,0) node[black,right] {$\vec F_\text{наблюдателя}$};



    \end{tikzpicture}
\caption{Наблюдатель толкает тележку}
\label{fig:p33}
\end{figure}

Наблюдатель~Н действует на тележку~Т вправо силой~$\vec F_\text{наблюдателя}$. Кроме того тележка притягивается к планете вниз силой~$\vec F_\text{планеты}$ и отталкивается от поверхности вверх силой~$\vec F_\text{поверхности}$; при этом $\vec F_\text{планеты}\hm+\vec F_\text{поверхности}\hm=0$. 

В этом случае\footnote{Под колесами гладкая поверхность.} на тележку действует результирующая сила, равная $\vec R_\text{т}\hm=\vec F_\text{планеты}+\vec F_\text{поверхности}+\vec F_\text{наблюдателя}=\vec F_\text{наблюдателя}\neq0$, и тележка приобретает ускорение~$\vec a_\text{т}$. Также наблюдатель заметил, что чем больше результирующая сила~$\vec R_\text{т}$~--- тем больше ускорение~$\vec a_\text{т}$, причем~$\vec R_\text{т}$ и~$\vec a_\text{т}$ сонаправлены. 

Связь \emph{результирующей силы} с ускорением дает второй закон Ньютона.

\begin{law}
\textbf{Второй закон Ньютона.} Результирующая сила, действующая на тело, равна произведению массы тела на его ускорение:
\begin{equation}
    \vec R= m\vec a.
\end{equation}
\end{law}

С другой стороны, наблюдатель ощущает действие на себя со стороны тележки, которая как бы <<мешает>> ему двигаться вправо (рис.\,\ref{fig:p34}).

\begin{figure}[H]
    \centering

    \begin{tikzpicture}[font=\small,            line cap=round,                             >={Stealth[inset=0pt,round,length=3mm, angle'=20]},vec/.style={>={Stealth[inset=0pt,round,length=3.5mm, angle'=20]}}]


\filldraw[lightgray,semitransparent] (0,0) rectangle (12,-0.3);
\draw (0,0) -- (12,0);

\begin{scope}[xshift=4cm]
    
%%motion1
\filldraw[lightgray,draw=black] (0,0.5) rectangle (4,1);
\node [above,black] at (2,1) {Т};
% \draw[->,Green,thick,vec] (3,1.2) --++ (1,0) node[black,above,midway] {$\vec a_\text{т}$};

%wheels
\draw[black,fill=black]
    (1,.3) circle(.3cm)
    (3,.3) circle(.3cm);
\draw[fill=gray]
    (1,.3) circle(.2cm)
    (3,.3) circle(.2cm);
\end{scope}



%%%%%%%%%%%%%%%%%%%%%%%%%%%%%%%%%%%%%%
% PERSON
  \def\W{3.6}  % ground width
  \def\D{0.2}  % ground depth
  \def\H{1.9}  % human height
  \def\h{0.9}  % mass height
  \def\w{1.0}  % mass width
  \def\mx{0.10*\W} % mass x coordinate

\begin{scope}[xshift=3.6cm]
    
  \coordinate (H) at (0,0.75*\H);
  \draw[thick,fill=white] (H) circle (0.3) node[xshift=-0.3cm,left] {$\text{Н}$};
  \draw[thick] (H)++(-120:0.3) coordinate (N) to[out=-115,in=40]++ (-130:0.36*\H) coordinate (P);
  \draw[thick] (N)++(-115:0.03) to[out=-60,in=170] (.4,\h);
  \draw[thick] (N)++(-115:0.03) to[out=-70,in=170] (.4,0.92*\h);
  \draw[thick] (P) to[out=-120,in=30,looseness=1.3] (-0.35*\W,0);
  \draw[thick] (P) to[out=-50,in=30,looseness=1.8] (-0.22*\W,0);

  %F
  % \node at (P) [circle,fill=NavyBlue,inner sep=0pt,minimum size=1mm,label=below:] {};
  % \draw[->,NavyBlue,thick,vec] (P) --++ (-3,0) node[black,above,midway] {$\vec F_\text{тележки}$};
      
\end{scope}  

%vec
\node at (6,.75) [circle,fill=red,inner sep=0pt,minimum size=1mm,label=below:] {};
\draw[->,red,thick,vec] (6,.75) --++ (3,0) node[black,right] {$\vec F_\text{наблюдателя}$};

%inter Ft
\path[name path=g] (6,.75) --++ (-3,0);
\path[thick,name path=h] (H)++(-120:0.3) coordinate (N) to[out=-115,in=40]++ (-130:0.36*\H) coordinate (P);
\node [name intersections={of=g and h, by=x1},circle,fill=NavyBlue,inner sep=0pt,minimum size=1mm,label=below:] at (x1) {};
\draw[->,NavyBlue,thick,vec] (x1) --++ (-3,0) node[black,above,midway] {$\vec F_\text{тележки}$};

    \end{tikzpicture}
\caption{Взаимодействие наблюдателя и тележки}
\label{fig:p34}
\end{figure}

Действительно, и тележка в рассматриваемом опыте действует на наблюдателя (сила~$\vec F_\text{тележки}$)~--- ведь если резко <<убрать>> тележку, то наблюдатель устремится туда, куда происходит толкание, то есть вправо. (Вообще любые силы носят \emph{взаимный} характер: если тело~A действует на тело~Б, то и тело~Б действует на тело~A.) 

\begin{law}
\textbf{Третий закон Ньютона.} Два тела действуют друг на друга с силами, равными по модулю
и противоположными по направлению (при этом обе силы лежат на одной прямой):
\begin{equation}
    \vec F_{1\to2}=-\vec F_{2\to1},
\end{equation}
где $\vec F_{1\to2}$~--- сила, действующая на второе тело со стороны первого; $\vec F_{2\to1}$~--- сила, действующая на первое тело со стороны второго.
\end{law}

Так, согласно третьему закону Ньютона: $F_\text{наблюдателя}= F_\text{тележки}$ (рис.\,\ref{fig:p34}). 







 
\end{document}